\newpage
\phantomsection
\label{prf:maximum-minimum-theorem}

\begin{remark*}[Return]
\hyperref[thm:maximum-minimum-theorem]{Return to Theorem}
\end{remark*}

\begin{theorem*}[Maximum-Minimum Theorem]
Every function continuous on a closed bounded interval $[a,b]$ attains its maximum and minimum on $[a,b]$.
\end{theorem*}

\begin{proof}
\textbf{Professional Standard Proof.}~\\
Let $I=[a,b]$ be a closed bounded interval and let $f : I \to \mathbb{R}$ be continuous on $I$. Since $f$ is continuous on the closed bounded interval $I$, the Boundedness Theorem implies that $f$ is bounded on $I$. Let $S:=\{f(x)\mid x\in I\}$. Since $S$ is bounded, its supremum and infimum exist. Let $s := \sup S$ and $i := \inf S$.

For each $n\in\mathbb{N}$, since $s-\frac{1}{n}$ is not an upper bound for $S$, there exists $x_n\in I$ such that $s-\frac{1}{n}<f(x_n)\le s$. The sequence $(x_n)\subseteq[a,b]$ is bounded, so by Bolzano-Weierstrass, there exists a subsequence $(x_{n_k})$ and a point $c\in I$ such that $x_{n_k}\to c$. By continuity of $f$ at $c$, we have $f(x_{n_k})\to f(c)$. Since $s-\frac{1}{n_k}<f(x_{n_k})\le s$ and $\frac{1}{n_k}\to 0$, the Squeeze Theorem implies $f(x_{n_k})\to s$. Therefore $f(c)=s$, so $f$ attains its supremum on $I$.

Similarly, for each $n\in\mathbb{N}$, there exists $y_n\in I$ such that $i\le f(y_n)<i+\frac{1}{n}$. By the same argument with Bolzano-Weierstrass and continuity, there exists $d\in I$ such that $f(d)=i$, so $f$ attains its infimum on $I$.
\end{proof}

\begin{proof}
\textbf{Detailed Learning Proof.}~\\

\textbf{Step 1.} Let $I=[a,b]$ be a closed bounded interval and let $f : I \to \mathbb{R}$ be continuous on $I$. Since $f$ is continuous on the closed bounded interval $I$, the Boundedness Theorem guarantees that $f$ is bounded on $I$. This allows us to consider the range $S:=\{f(x)\mid x\in I\}$, which is a bounded subset of $\mathbb{R}$.

\textbf{Step 2.} Since $S$ is bounded above and below, the completeness of $\mathbb{R}$ ensures that $\sup S$ and $\inf S$ exist. Let $s := \sup S$ and $i := \inf S$. These represent the least upper bound and greatest lower bound of the function values, respectively.

\textbf{Step 3.} For each $n\in\mathbb{N}$, since $s$ is the supremum, $s-\frac{1}{n}$ cannot be an upper bound for $S$. Therefore, there exists $x_n\in I$ such that $s-\frac{1}{n}<f(x_n)\le s$. This constructs a sequence $(x_n)$ in $I$ whose function values approach the supremum.

\textbf{Step 4.} The sequence $(x_n)\subseteq[a,b]$ is bounded since $[a,b]$ is bounded. By the Bolzano-Weierstrass Theorem, there exists a subsequence $(x_{n_k})$ and a point $c\in[a,b]$ such that $x_{n_k}\to c$. The closed interval $[a,b]$ contains all its limit points, so $c\in I$.

\textbf{Step 5.} Since $f$ is continuous at $c$, the convergence $x_{n_k}\to c$ implies $f(x_{n_k})\to f(c)$. Simultaneously, from the construction in Step 3, we have $s-\frac{1}{n_k}<f(x_{n_k})\le s$. Since $\frac{1}{n_k}\to 0$, the Squeeze Theorem forces $f(x_{n_k})\to s$. Therefore $f(c)=s$, proving that $f$ attains its supremum at $c$.

\textbf{Step 6.} For the infimum, construct a sequence $(y_n)$ in $I$ such that $i\le f(y_n)<i+\frac{1}{n}$ for each $n\in\mathbb{N}$. Apply the same subsequence extraction and continuity argument to find $d\in I$ with $f(d)=i$, proving that $f$ attains its infimum.
\end{proof}

\begin{remark*}[Proof structure]
This proof employs a direct constructive approach combined with sequential compactness. The strategy first establishes boundedness of the range using the Boundedness Theorem, then explicitly constructs sequences whose function values approach the supremum and infimum. The key insight is using the Bolzano-Weierstrass Theorem to extract convergent subsequences from the bounded domain, then applying continuity to show the limit points achieve the extreme values. The compactness of the closed bounded interval $[a,b]$ is essential for guaranteeing convergent subsequences.
\end{remark*}

\begin{remark*}[Dependencies]~\\
\begin{itemize}
  \item \hyperref[def:absolute-extrema]{Definition of absolute maximum and absolute minimum}
  \item \hyperref[thm:boundedness-theorem]{Boundedness Theorem}
  \item \hyperref[thm:completeness-real-numbers]{Completeness of $\mathbb{R}$}
  \item \hyperref[thm:bolzano-weierstrass]{Bolzano-Weierstrass Theorem}
  \item \hyperref[thm:squeeze-theorem]{Squeeze Theorem}
\end{itemize}
\end{remark*}
